\section{Parité}
Les tests que nous avons faits jusqu'à présent dans les instructions
``if'' concernaient l'orientation du \robot.  En \python\ il est
possible de faire des tests concernant les valeurs des paramètres et des
compteurs de boucles ``for''.  Il peut être utile de faire un test sur le compteur
dans le corps d'une boucle pour exécuter des actions différentes à
chaque itération.  Prenons l'exemple du dessin de la frise en forme de
double peigne ci-dessous~:

 {\raggedright
  \includevspacePostScript[1.5cm]{FriseDoublePeigne}{1.0}{1.0}
   }

Les dents du peigne étant dessinées de gauche à droite, les
$1^{\mathrm{\grave{e}re}}$, $3^{\mathrm{\grave{e}me}}$,
$5^{\mathrm{\grave{e}me}}$, \ldots\ sont dirigées vers le bas alors
que les $2^{\mathrm{\grave{e}me}}$, $4^{\mathrm{\grave{e}me}}$,
$6^{\mathrm{\grave{e}me}}$, \ldots\ sont dirigées vers le haut.  Dans
le programme de dessin on utilise un compteur de boucle {\idt i} qui
vaut 1 quand on dessine la $1^{\mathrm{\grave{e}re}}$ dent, 2 quand on
dessine la $2^{\mathrm{\grave{e}me}}$, \ldots\ Par conséquent quand la
valeur du compteur de boucle est {\em impaire} (c'est-à-dire 1, 3, 5,
\ldots) il faut dessiner une dent vers le bas tandis que si la valeur
du compteur est {\em paire} (c'est-à-dire 2, 4, 6, \ldots) il faut
dessiner une dent vers le haut.  \index{Parité}Un nombre est impair
quand le reste de sa division par 2 est égal à 1. Pour tester que la
valeur du compteur {\idt i} est impaire on utilise donc l'expression
booléenne {\idt i \% 2 == 1} qui vaut {\mcl True} (vrai) quand le reste
{\idt i \% 2} de la division entière de {\idt i} par 2 est égal à 1 et
{\mcl False} (faux) quand ce reste est égal à 0, c'est-à-dire quand
{\idt i} est pair.  On remarquera que le test d'égalité s'écrit avec deux
signes égal {\idt ==}, le signe {\idt =} seul étant réservé à la
déclaration des constantes comme {\idt nombreDeDents = 30}.

{
\parindent \parindentprog
\medskip

\carrenoir
\vspace{\neglskip}\mbox{\blc {\mcl from}{\idt \ LibRobot\ }{\mcl import}{\idt \ *}}\par
\vspace{0.3\baselineskip}
\vspace{\neglskip}\mbox{\blc {\cmt \#\ Dessiner\ une\ frise\ en\ forme\ de\ double\ peigne}}\par
\vspace{\neglskip}\mbox{\blc {\idt nombreDeDents\ =\ 30}}\par
\vspace{0.3\baselineskip}
\vspace{\neglskip}\mbox{\blc {\idt ag();\ pe()}}\par
\vspace{\neglskip}\mbox{\blc {\mcl for}{\idt \ i\ }{\mcl in}{\idt \ range(1,\ nombreDeDents\ +\ 1):}}\par
\vspace{\neglskip}\mbox{\blc \blk\blk{\cmt \#\ Trait\ horizontal}}\par
\vspace{\neglskip}\mbox{\blc \blk\blk{\idt av()}}\par
\vspace{\neglskip}\mbox{\blc \blk\blk{\mcl if}{\idt \ i\ \%\ 2\ ==\ 1:}}\par
\vspace{\neglskip}\mbox{\blc \blk\blk\blk\blk{\cmt \#\ Dent\ vers\ le\ bas}}\par
\vspace{\neglskip}\mbox{\blc \blk\blk\blk\blk{\idt pqtd();\ av();\ pdt();\ av();\ pqtd()}}\par
\vspace{\neglskip}\mbox{\blc \blk\blk{\cmt \#\ Dent\ vers\ le\ haut}}\par
\vspace{\neglskip}\mbox{\blc \blk\blk{\mcl else}{\idt :}}\par
\vspace{\neglskip}\mbox{\blc \blk\blk\blk\blk{\idt pqtg();\ av();\ pdt();\ av();\ pqtg()}}\par
\vspace{\neglskip}\mbox{\blc {\idt st()}}\par

\medskip }


\section{Comparaison d'entiers}

Si {\idt i} est un compteur de boucle ou un paramètre entier de procédure,
l'expression booléenne {\idt i > 10} est {\mcl True} (vraie)
si la valeur de {\idt i} est strictement plus grande que 10 et {\mcl
False} (fausse) dans le cas contraire c'est-à-dire si la valeur de {\idt
i} est inférieure ou égale à 10.  Les symboles mathématiques $\leq$
(inférieur ou égal) et $\geq$ (supérieur ou égal) se notent {\idt <=} et
{\idt >=} en \python\ (car les symboles mathématiques $\leq$ et $\geq$ ne
se trouvent pas sur les claviers d'ordinateurs alors que {\idt <}, {\idt >}
et {\idt =} sont toujours présents).  Par exemple les expressions
booléennes {\idt 10 <= 10} et {\idt 10 <= 20} valent {\mcl True}
(vrai) tandis que {\idt 10 <= 9} vaut {\mcl False} (faux).  En
\python, comme nous l'avons vu, on note l'égalité avec deux signes égal {\idt ==}
puisque le symbole mathématique usuel {\idt =} sert à déclarer les constantes.
Le symbole mathématique $\neq$ (différent) se note {\ttfamily\idt !=}.
Par exemple {\idt 10 == 10} et {\ttfamily\idt 10 != 20} valent {\mcl True}
(vrai) tandis que {\ttfamily\idt 10 != 10} et {\idt 10 == 20} valent {\mcl
False} (faux)\footnote{\ Si l'on écrit par erreur {\mcl if}\ {\idt i = 10:}
au lieu de {\mcl if}\ {\idt i == 10:}, \python\ refuse le programme avec
le message {\idt SyntaxError: invalid syntax. Maybe you meant '==' or ':='
instead of '='?}}.
On peut utiliser les expressions booléennes de
comparaison d'entiers ou de valeurs de paramètres entiers dans le test
de l'instruction ``if''.  Par exemple la commande {\idt avf(n)} qui
fait avancer le \robot\ {\idt n} fois si la valeur de {\idt n} est
positive, le fait reculer si la valeur de {\idt n} est négative et qui
ne change pas sa position quand la valeur de {\idt n} est nulle peut
se définir comme suit~:

\input python/DefinitionAvf.tex

\noindent On remarquera que le paramètre effectif négatif de l'appel
{\idt avf(-5)} s'écrit simplement avec son signe, entre les parenthèses de
l'appel. Pour reculer, le \robot\ fait demi-tour, avance {\idt -n} fois
(qui est positif) et refait demi-tour.

De manière plus générale, on peut utiliser les opérateurs {\idt <} (strictement inférieur), {\idt <=} (inférieur ou égal),
{\idt ==} (égal), {\ttfamily\idt !=} (différent), {\idt >=} (supérieur ou égal) et {\idt >} (strictement supérieur) pour comparer les valeurs
de deux expressions entières comme par exemple {\idt i <= l} dans le programme suivant qui trace le serpentin
ci-dessous. La commande {\idt lg} accepte aussi bien un nombre entier, comme
dans {\idt lg(20)} ou {\idt lg(i)}, qu'un nombre avec une virgule, écrite
avec un point comme dans {\idt lg(20.5)} (nous y reviendrons
au chapitre 19, consacré aux expressions rationnelles)~:

 \includevspacePostScript[1.5cm]{Serpentin}{0.95}{0.95}
 \input python/Serpentin.tex

\section{Négation, conjonction et disjonction}

La {\em négation} \index{Négation} d'une expression booléenne $E$ s'écrit {\mcl not} $E$, le mot anglais {\mcl not} signifiant `non'. Si la
valeur de $E$ est {\mcl True} alors celle de {\mcl not} $E$ est {\mcl False} tandis que si la valeur de $E$ est {\mcl False} alors
celle de {\mcl not} $E$ est {\mcl True}. Par exemple  {\mcl not} {\idt (5 \% 2 == 1)} vaut {\mcl False} puisque
{\idt 5 \% 2 == 1} vaut {\mcl True}, {\idt 5} étant impair. Les parenthèses ne sont
pas indispensables mais rendent l'expression plus lisible (nous verrons pourquoi dans l'Addenda). Ceci peut se résumer par la table de vérité suivante~:

\begin{center}
\begin{tabular}{|c|c|} \hline
$E$ & {\mcl not} $E$\\ \hline \hline
{\mcl True} & {\mcl False}\\ \hline
{\mcl False} & {\mcl True}\\ \hline
\end{tabular}
\end{center}

La {\em conjonction} \index{Conjonction} de deux expressions booléennes $E_{1}$ et $E_{2}$ s'écrit $E_{1}$ {\mcl and} $E_{2}$,
ce qui se lit `$E_{1}$ et $E_{2}$'. La valeur de $E_{1}$ {\mcl and} $E_{2}$ est {\mcl True} (vrai) si et
seulement si la valeur de $E_{1}$ et celle de $E_{2}$ sont {\mcl True} (vrai). Si l'une des expressions
booléennes $E_{1}$ ou
$E_{2}$ vaut {\mcl False} (faux) alors la conjonction $E_{1}$ {\mcl and} $E_{2}$ vaut également {\mcl False} (faux).
Par exemple  {\idt 5 >= 0} {\mcl and} {\idt 5 \% 2 == 1} vaut {\mcl True} puisque
{\idt 5} est positif et impair. Ceci peut se résumer par la table de vérité suivante~:

\begin{center}
\begin{tabular}{|c|c|c|} \hline
$E_{1}$ & $E_{2}$ & $E_{1}$ {\mcl and} $E_{2}$\\ \hline \hline
{\mcl True} & {\mcl True} & {\mcl True}\\ \hline
{\mcl True} & {\mcl False} & {\mcl False}\\ \hline
{\mcl False} & {\mcl True} & {\mcl False}\\ \hline
{\mcl False} & {\mcl False} & {\mcl False}\\ \hline
\end{tabular}
\end{center}

La {\em disjonction} \index{Disjonction} de deux expressions booléennes $E_{1}$ et $E_{2}$ s'écrit $E_{1}$ {\mcl or} $E_{2}$,
ce qui se lit `$E_{1}$ ou $E_{2}$'. La valeur de $E_{1}$ {\mcl or} $E_{2}$ est {\mcl True} (vrai) si et
seulement si la valeur de $E_{1}$ ou celle de $E_{2}$ est {\mcl True} (vrai). Si les valeurs de $E_{1}$ et de
$E_{2}$ sont  {\mcl False} (faux) alors la disjonction $E_{1}$ {\mcl or} $E_{2}$ vaut également {\mcl False} (faux).
Par exemple  {\idt 4 >= 0} {\mcl or} {\idt 4 \% 2 == 1} vaut {\mcl True} puisque
{\idt 4} est positif et bien que {\idt 4} soit pair. Ceci peut se résumer par la table de vérité suivante~:

\begin{center}
\begin{tabular}{|c|c|c|} \hline
$E_{1}$ & $E_{2}$ & $E_{1}$ {\mcl or} $E_{2}$\\ \hline \hline
{\mcl True} & {\mcl True} & {\mcl True}\\ \hline
{\mcl True} & {\mcl False} & {\mcl True}\\ \hline
{\mcl False} & {\mcl True} & {\mcl True}\\ \hline
{\mcl False} & {\mcl False} & {\mcl False}\\ \hline
\end{tabular}
\end{center}

La conjonction et la disjonction peuvent se réaliser simplement avec un circuit électrique comme on le voit à la figure
\ref{fig:andor}
ci-dessous
%(page \pageref{fig:andor})
avec la con\-ven\-tion que {\mcl True} représente l'interrupteur fermé
et l'ampoule allumée (le courant passe) et {\mcl False} représente l'interrupteur ouvert
et l'ampoule éteinte (le courant ne passe pas).

\begin{figure}[htbp]
 \includevspacePostScript[4.3cm]{CircuitsAndOr}{0.95}{0.95}
 \caption{Circuits électriques réalisant la conjonction et\label{fig:andor}}
\centering{la disjonction logiques}
\end{figure}

Une disjonction est utilisée dans le programme suivant pour tester si le \robot\ est orienté au nord ou à
l'est~:

\vspace{0.3cm}
 {\raggedright
 \includevspacePostScript[0.5cm]{FriseEnOr}{1.0}{1.0}
{
\parindent \parindentprog
\medskip

\carrenoir
\vspace{\neglskip}\mbox{\blc {\mcl from}{\idt \ LibRobot\ }{\mcl import}{\idt \ *}}\par
\vspace{0.3\baselineskip}
\vspace{\neglskip}\mbox{\blc {\cmt \#\ Dessiner\ une\ frise\ en\ or}}\par
\vspace{\neglskip}\mbox{\blc {\idt longueur\ =\ 31}}\par
\vspace{0.3\baselineskip}
\vspace{\neglskip}\mbox{\blc {\idt ag()}}\par
\vspace{\neglskip}\mbox{\blc {\mcl for}{\idt \ i\ }{\mcl in}{\idt \ range(longueur):}}\par
\vspace{\neglskip}\mbox{\blc \blk\blk{\mcl if}{\idt \ tn()\ }{\mcl or}{\idt \ te():}}\par
\vspace{\neglskip}\mbox{\blc \blk\blk\blk\blk{\idt vd()}}\par
\vspace{\neglskip}\mbox{\blc \blk\blk{\mcl else}{\idt :}}\par
\vspace{\neglskip}\mbox{\blc \blk\blk\blk\blk{\idt avf({-}1);\ vg();\ pqtg()}}\par
\vspace{\neglskip}\mbox{\blc {\idt st()}}\par

\medskip }

  }

\section{Expressions booléennes}

La table ci-dessous résume les règles d'écriture et les règles de calcul de la valeur
des expressions booléennes en \python~:

\index{Expression booléenne} \index{Test}

\begin{center}
\protect\frenchspacing
\begin{tabular}{@{}|l|p{6.8cm}|@{}} \hline
\multicolumn{2}{|l|}{\em Expressions booléennes~:} \\ \hline \hline
{\mcl True} & Valeur `vrai'.\\
{\mcl False} & Valeur `faux'.\\ \hline
{\small\rm Expr. entière$_{1}$} {\idt <=} {\small\rm Expr. entière$_{2}$}& {\mcl True} si et seulement si la valeur de l'{\small\rm Expr. entière$_{1}$} est
                                 inférieure ou égale à celle de l'{\small\rm Expr. entière$_{2}$}.\\
{\small\rm Expr. entière$_{1}$} {\idt <} {\small\rm Expr. entière$_{2}$}& {\mcl True} si et seu\-le\-ment si la va\-leur de l'{\small\rm Expr. entière$_{1}$} est
                                 strictement inférieure à celle de l'{\small\rm Expr. entière$_{2}$}.\\
{\small\rm Expr. entière$_{1}$} {\idt ==} {\small\rm Expr. entière$_{2}$}& {\mcl True} si et seu\-le\-ment si la valeur de l'{\small\rm Expr. entière$_{1}$} est
                                 égale à celle de l'{\small\rm Expr. entière$_{2}$}.\\
{\small\rm Expr. entière$_{1}$} {\ttfamily\idt !=} {\small\rm Expr. entière$_{2}$}& {\mcl True} si et seu\-le\-ment si la valeur de l'{\small\rm Expr. entière$_{1}$} est
                                 différente de celle de l'{\small\rm Expr. entière$_{2}$}.\\ \hline
\end{tabular}
\end{center}
\begin{center}
\protect\frenchspacing
\begin{tabular}{@{}|l|p{6.8cm}|@{}} \hline
{\small\rm Expr. entière$_{1}$} {\idt >} {\small\rm Expr. entière$_{2}$}& {\mcl True} si et seu\-le\-ment si la va\-leur de l'{\small\rm Expr. entière$_{1}$} est
                                 strictement supérieure à celle de l'{\small\rm Expr. entière$_{2}$}.\\
{\small\rm Expr. entière$_{1}$} {\idt >=} {\small\rm Expr. entière$_{2}$}& {\mcl True} si et seulement si la valeur de l'{\small\rm Expr. entière$_{1}$} est
                                supérieure ou égale à celle de l'{\small\rm Expr. entière$_{2}$}.\\ \hline
                                \end{tabular}
\end{center}
\begin{center}
\protect\frenchspacing
\begin{tabular}{@{}|l|p{6.8cm}|@{}} \hline
{\mcl not} {\small\rm Expr. bool.} & {\mcl True} si et seulement si la valeur de l'{\small\rm Expr. bool.} est
                                égale à {\mcl False}.\\ \hline
{\small\rm Expr. bool.$_{1}$} {\mcl and} {\small\rm Expr. bool.$_{2}$}& {\mcl True} si et seulement si les valeurs de l'{\small\rm Expr. bool.$_{1}$} et
                                 de l'{\small\rm Expr. bool.$_{2}$} sont égales à {\mcl True}.\\
{\small\rm Expr. bool.$_{1}$} {\mcl or} {\small\rm Expr. bool.$_{2}$}& {\mcl True} si et seulement si la valeur de l'{\small\rm Expr. bool.$_{1}$}
                                 ou celle de l'{\small\rm Expr. bool.$_{2}$} est égale à {\mcl True}.\\ \hline
\end{tabular}
\end{center}

\index{Fonction booléenne}On peut donner un nom à une expression booléenne pour
rendre les programmes plus lisibles. Par exemple~:

{
\parindent \parindentprog
\medskip

\vspace{\neglskip}\mbox{\blc {\cmt \#\ n\ est-il\ impair\ ?}}\par
\vspace{\neglskip}\mbox{\blc {\mcl def}{\idt \ estImpair(n):}}\par
\vspace{\neglskip}\mbox{\blc \blk\blk{\mcl return}{\idt \ n\ \%\ 2\ ==\ 1}}\par
\vspace{0.3\baselineskip}
\vspace{\neglskip}\mbox{\blc {\cmt \#\ n\ est-il\ pair\ ?}}\par
\vspace{\neglskip}\mbox{\blc {\mcl def}{\idt \ estPair(n):}}\par
\vspace{\neglskip}\mbox{\blc \blk\blk{\mcl return}{\idt \ n\ \%\ 2\ ==\ 0}}\par

\medskip }

\noindent Ces déclarations ressemblent à des déclarations de procédures avec paramètres
mais leur corps n'est pas une suite d'instructions pour le
\robot~: l'instruction {\mcl return} (`retourner' en anglais) est suivie d'une
expression booléenne dont la valeur est calculée à chaque appel et retournée comme
résultat de l'appel.
On dit que {\idt estImpair} et {\idt estPair} sont des {\em fonctions booléennes} car leur évaluation retourne toujours un booléen (une valeur de type {\idt bool}, c'est-à-dire {\mcl True} ou {\mcl False}).
Par exemple {\idt estImpair(5)} vaut {\mcl True} et {\idt estPair(5)} vaut
{\mcl False}. On peut les utiliser dans un test comme {\mcl if}\ {\idt estImpair(i):}
\ldots\ ou dans une expression booléenne comme {\mcl not}\ {\idt estImpair(i)}. Les tests
d'orientation {\idt tn()}, {\idt te()}, \ldots\ du \robot\ sont également des fonctions
booléennes, déclarées dans la bibliothèque {\idt LibRobot}.

\noindent \rule{137mm}{0.1mm} %---------------------------------------------
\begin{exercices}{Écrire des programmes \python\  utilisant des tests de parité pour réaliser
les dessins suivants~:

\tiret  la serpentine pyramidale et la grille de la figure \ref{fig:serpyrgricr} ci-dessous$\:$;
%(page \pageref{fig:serpyrgricr})$\:$;

\tiret   le pavage octogonal et le pavage d'hexagones entrelacés de la figure \ref{fig:pavocthex} (page \pageref{fig:pavocthex}).

\noindent $\Box$}\end{exercices}
\begin{figure}[htbp]
  \includevspacePostScript[6.0cm]{SerpPyr+GrilleDeCr}{0.95}{0.95}
 \caption{Serpentine pyramidale et grille de croix}\label{fig:serpyrgricr}
\end{figure}
\begin{figure}[htbp]
\includevspacePostScript[5.4cm]{PavageOctoHexa}{0.95}{0.95}
 \caption{Pavages d'octogones juxtaposés et d'hexagones entrelacés}\label{fig:pavocthex}
\end{figure}
\begin{exercices}{Écrire des programmes \python\  utilisant des expressions boolé\-en\-nes de comparaison de
valeurs entières pour réaliser les dessins suivants (le deuxième est difficile)~:

\tiret  les chiffres digitaux de la figure \ref{fig:alphabet} (page \pageref{fig:alphabet}) où le chiffre à dessiner
est passé en paramètre de la procédure de dessin$\:$;

\tiret  le pavage de roues dentées de la figure \ref{fig:pavrd} ci-dessous$\:$;

\tiret  le tapis égyptien de la figure \ref{fig:tapegyp} (page \pageref{fig:tapegyp}).

\noindent $\Box$}\end{exercices}
\begin{figure}[htbp]
 \includevspacePostScript[5.7cm]{PavageRoueDentee}{0.95}{0.95}
 \caption{Pavage de roues dentées}\label{fig:pavrd}
\end{figure}
\begin{figure}[htbp]
 \includevspacePostScript[7.2cm]{TapisEgyptien}{0.95}{0.95}
 \caption{Tapis égyptien\label{fig:tapegyp}}
\end{figure}
\begin{exercices}{Écrire des programmes \python\  utilisant des négations, conjonctions ou disjonctions d'expressions
booléen\-nes pour réaliser les dessins suivants (le deuxième est difficile)~:

\tiret  le carré zébré de la figure \ref{fig:carzeb} (page \pageref{fig:carzeb})$\:$;

\tiret  la mosaïque de l'Alhambra représentée à la figure \ref{fig:mosaic6} (page \pageref{fig:mosaic6}).

\noindent $\Box$}\end{exercices}
\begin{figure}[htbp]
 \includevspacePostScript[5.8cm]{CarreZebree}{0.95}{0.95}
 \caption{Carré zébré}\label{fig:carzeb}
\end{figure}
\begin{figure}[htbp]
 \includevspacePostScript[5.7cm]{MosaiqueAlhambra6}{0.95}{0.95}
 \caption{Mosaïque de l'Alhambra}\label{fig:mosaic6}
\end{figure}
%%%%%%%%%%%%%%%%%%%%%%%%%%%%%%%%%%%%%%%%%%%%%%%%%%%%%%%%%%
\begin{corrige}{
Nous donnons un programme, les autres se trouvant dans le dossier
{\idt livre/\allowbreak 16\_expressions\_booleennes}~:
\input python/SerpentinePyramidale.tex
\noindent $\Box$}\end{corrige}
\begin{corrige}{
Nous donnons un programme, le troisième se trouvant dans le dossier
{\idt livre/\allowbreak 16\_expressions\_booleennes}~:
\input python/PavageRouesDentees.tex
\noindent $\Box$}\end{corrige}
\begin{corrige}{
Nous donnons un programme, le premier se trouvant dans le dossier
{\idt livre/\allowbreak 16\_expressions\_booleennes}~:
\input python/MosaiqueAlhambra6.tex
\noindent $\Box$}\end{corrige}

%%%%%%%%%%%%%%%%%%%%%%%%%%%%%%%%%%%%%%%%%%%%%%%%%%%%%%%%%%%%%%%%%%%%%%%%%%%
\noindent \framebox[\textwidth]{
\begin{minipage}{\myparboxwidth}
\setlength{\parindent}{2em}
\small \sf
\vspace{\abovedisplayskip}
\section*{Addenda}

 En \python\ les priorités des opérateurs entiers et booléens sont définies comme suit~:

 \vspace{\baselineskip}

 \index{Priorité des opérateurs}
\begin{center}
\begin{tabular}{|l|l|l|} \hline
\multicolumn{3}{|l|}{\em Priorité des opérateurs entiers et booléens~:} \\ \hline \hline
{\idt f(x)} (appel) & appels de fonctions & priorité la plus forte \\ \hline
$-$ (unaire) & opposé & \\ \hline
$\ast$, {\idt //}, {\idt \%} & opérateurs multiplicatifs & \\ \hline
$+$, $-$ (binaires) & opérateurs additifs & \\ \hline
{\idt <}, {\idt <=}, {\idt ==}, {\ttfamily\idt !=}, {\idt >=}, {\idt >}  & opérateurs de relation & \\ \hline
{\mcl not} & négation & \\ \hline
{\mcl and} & conjonction & \\ \hline
{\mcl or} & disjonction & priorité la plus faible \\ \hline
\end{tabular}
\end{center}

\vspace{\baselineskip}

 \noindent Par exemple l'expression booléenne `{\mcl True\ and\ False\ or\ False}' s'é\-va\-lue com\-me
 `{\idt (}{\mcl True\ and\ False}{\idt )}\ {\mcl or\ False}' et est donc égale à {\mcl False}.

\vspace{\abovedisplayskip}
\end{minipage}
}
%%%%%%%%%%%%%%%%%%%%%%%%%%%%%%%%%%%%%%%%%%%%%%%%%%%%%%%%%%%%%%%%%%%%%%%%%%%%

%%%%%%%%%%%%%%%%%%%%%%%%%%%%%%%%%%%%%%%%%%%%%%%%%%%%%%%%%%%%%%%%%%%%%%%%%%%
\noindent \framebox[\textwidth]{
\begin{minipage}{\myparboxwidth}
\setlength{\parindent}{2em}
\small \sf
\vspace{\abovedisplayskip}
\indent Ces priorités sont proches des conventions mathématiques usuelles. Par
exemple `{\idt 0\ <=\ i}\ {\mcl and}\ {\idt i\ <\ 100}' s'évalue comme
`{\idt (0\ <=\ i)}\ {\mcl and}\ {\idt (i\ <\ 100)}', ce qui est bien ce que l'on veut dire. Il n'en était
pas de même dans des langages plus anciens comme \pascal, où cette expression était
incorrecte et où il fallait écrire toutes les parenthèses. \python\ permet même
d'écrire `{\idt 0\ <=\ i\ <\ 100}' comme en mathématiques.

Comme {\mcl not} est moins prioritaire que les opérateurs de relation,
`{\mcl not}\ {\idt i\ \%\ 2\ ==\ 1}' s'évalue comme `{\mcl not}\ {\idt (i\ \%\ 2\ ==\ 1)}',
ce qui est bien ce que l'on veut dire. Il n'en est pas de même dans d'autres langages
comme \pascal\ ou \ocaml, où {\mcl not} s'applique à ce qui le suit immédiatement,
à la manière d'un appel de fonction, et où il faut donc écrire les parenthèses.
Dans le doute, on peut toujours ajouter des parenthèses.

Enfin, \python\ évalue les conjonctions et les disjonctions de gauche à droite et s'arrête dès
que le résultat est connu~: dans $E_{1}$ {\mcl and} $E_{2}$, si $E_{1}$ vaut {\mcl False}
alors $E_{2}$ n'est pas calculée puisque la conjonction vaut forcément {\mcl False}$\:$; de
même dans $E_{1}$ {\mcl or} $E_{2}$, si $E_{1}$ vaut {\mcl True} alors $E_{2}$ n'est pas
calculée.

Le terme ``booléen'' ou ``booléien'' ({\idt bool} en \python) vient de Boole\footnote{\ \href{https://fr.wikipedia.org/wiki/George_Boole}{George Boole}, mathématicien britannique, 1815--1864.},
 créateur de la logique mathématique moderne.

\vspace{\abovedisplayskip}
\end{minipage}
}
%%%%%%%%%%%%%%%%%%%%%%%%%%%%%%%%%%%%%%%%%%%%%%%%%%%%%%%%%%%%%%%%%%%%%%%%%%%%
\newpage                 % Si on termine le chapitre par une page blanche
\thispagestyle{empty}
